\chapter{Parent-origin центрального гамильтониана дополнительных масс}
\label{ch:v8-extra-mass-central-hamiltonian-parent-origin}

\section{Ортогональный базис двух щелей}

На центре блоков $(Y,u,d)$ введём
\begin{equation}
 A=P_u-P_d,\qquad B=P_u+P_d-3P_Y.
 \label{eq:v8-extra-mass-origin-basis}
\end{equation}
С учётом рангов $(4,6,6)$ эти операторы бесследовы и ортогональны:
\begin{equation}
 \Tr A^2=12,\qquad \Tr B^2=48,\qquad \Tr AB=0.
 \label{eq:v8-extra-mass-origin-gram}
\end{equation}
Для $h^0=aA+bB$ две щели равны
\begin{equation}
 \Delta_u=a+4b,\qquad \Delta_d=-a+4b,
 \label{eq:v8-extra-mass-origin-gap-map}
\end{equation}
а обратная карта имеет вид
\begin{equation}
 a=\frac{\Delta_u-\Delta_d}{2},\qquad
 b=\frac{\Delta_u+\Delta_d}{8}.
 \label{eq:v8-extra-mass-origin-gap-inverse}
\end{equation}

\section{Обычные спектральные моменты дают только жёсткость}

Первые два нетривиальных момента равны
\begin{equation}
 \Tr(h^0)^2=12a^2+48b^2,
 \qquad
 \Tr(h^0)^3=36a^2b-96b^3.
 \label{eq:v8-extra-mass-origin-moments}
\end{equation}
На скалярном фоне $\varepsilon I$ линейный отклик любого полиномиального
момента вдоль $A$ и $B$ исчезает из-за бесследовости. Для квадратичного
момента
\begin{equation}
 \nabla_{a,b}\Tr(\varepsilon I+h^0)^2\big|_0=0,
 \qquad
 \Hess_{a,b}=\operatorname{diag}(24,96).
 \label{eq:v8-extra-mass-origin-quadratic-stiffness}
\end{equation}
Следовательно, незавешенный spectral trace предоставляет положительную
жёсткость, но выбирает $(a,b)=(0,0)$. Новые скаляры и их связи могут
возникать после расширения конечной тройки, однако их Dirac-блоки и
коэффициенты являются частью расширенных данных
\cite{v8-extra-mass-origin-stephan,v8-extra-mass-origin-krajewski}.

\section{Gauge-Casimir near miss}

Квадратичные значения $Y^2,C_2,C_3$ дают условную карту коэффициентов
$c=(c_1,c_2,c_3)$ в щели
\begin{equation}
 \begin{pmatrix}\Delta_u\\\Delta_d\end{pmatrix}
 =\begin{pmatrix}
 91/36&-3/4&4/3\\
 7/36&-3/4&4/3
 \end{pmatrix}
 \begin{pmatrix}c_1\\c_2\\c_3\end{pmatrix}.
 \label{eq:v8-extra-mass-origin-gauge-map}
\end{equation}
Матрица имеет ранг два и одномерное ядро. При этом
\begin{equation}
 \Delta_u-\Delta_d=\frac73c_1.
 \label{eq:v8-extra-mass-origin-hypercharge-splitting}
\end{equation}
Поэтому равные щели требуют $c_1=0$, хотя оба цветных канала имеют разные
ненулевые гиперзаряды. Gauge-Casimir-операторы охватывают плоскость формы,
но текущий spectral action использует их в gauge-кинетике, а не как
центральный endpoint-гамильтониан. Перенос коэффициентов без нового
типизированного члена был бы смешением функционалов.

Grading на трёх edge-секторах скалярен и исчезает после quotient. Радиус
когерентности $\Tr(B_{
m coh}B_{
m coh}^*)=3$ также является скаляром:
\begin{equation}
 \Tr(3IA)=\Tr(3IB)=0.
 \label{eq:v8-extra-mass-origin-coherence-zero}
\end{equation}
Он не создаёт ни одной щели без заранее выбранного портального направления.

\section{Минимальный недостающий источник}

Условный двухкомпонентный source-parent равен
\begin{equation}
 V(a,b)=\frac12\Tr(h^0)^2-j_Aa-j_Bb
 =6a^2+24b^2-j_Aa-j_Bb.
 \label{eq:v8-extra-mass-origin-source-preview}
\end{equation}
Его единственная стационарная точка есть
\begin{equation}
 a_*=\frac{j_A}{12},\qquad b_*=\frac{j_B}{48}.
 \label{eq:v8-extra-mass-origin-source-stationary}
\end{equation}
Для общей точки симплекса нужны две независимые source-компоненты. Старое
действие не зависит от нового центра и даёт $(j_A,j_B)=(0,0)$.

Итоговые реестры:
\begin{equation}
 N_{\rm conditional\ shape}=5/5,
 \qquad N_{\rm parent\ origin}=0/7.
 \label{eq:v8-extra-mass-origin-ledgers}
\end{equation}

\begin{theorem}[Жёсткость без ненулевого источника]
Трассовый квадрат канонически задаёт положительную метрику на двумерной
плоскости центральных щелей, а gauge-Casimir-данные условно охватывают оба
направления. Однако обычные моменты на скалярном фоне имеют нулевой линейный
отклик, grading и coherence-радиус скалярны, а gauge-коэффициенты не входят
в endpoint-гамильтониан. Поэтому ненулевая пара щелей не выведена.
\end{theorem}

\begin{nogocriterion}[Запрет импорта gauge-коэффициентов как энергий]
Нельзя считать коэффициенты gauge-кинетики энергиями центральных edge-
уровней без единого оператора, в котором оба появления получены одним
следом. Нельзя также умножать скалярный coherence-конденсат на $A$ или $B$
без происхождения самого портального направления.
\end{nogocriterion}

\begin{protocol}\begin{enumerate}
\item базис щелей: \textbf{$A,B$};
\item Gram-метрика: \textbf{$\operatorname{diag}(12,48)$};
\item обычный след создаёт жёсткость: \textbf{да};
\item обычный след создаёт ненулевой источник: \textbf{нет};
\item gauge-Casimir map: \textbf{ранг $2$};
\item равные щели требуют: \textbf{$c_1=0$};
\item coherence-скаляр создаёт щель: \textbf{нет};
\item условная форма: \textbf{$5/5$};
\item parent-origin: \textbf{$0/7$};
\item минимально недостающих источников: \textbf{два};
\item следующий гейт: \textbf{архитектура двухкомпонентного source-parent}.
\end{enumerate}\end{protocol}

Машинный сертификат сохранён в
\path{s2t_v8_baryon_c0_singlet_triplet_central_gap_isotypic_channel_extra_edge_mass_central_hamiltonian_parent_action_origin_gate_results.json}.

\begin{thebibliography}{9}
\bibitem{v8-extra-mass-origin-stephan} C.~A.~Stephan,
\emph{New Scalar Fields in Noncommutative Geometry}, arXiv:0901.4676.
\bibitem{v8-extra-mass-origin-krajewski} T.~Krajewski,
\emph{Classification of Finite Spectral Triples}, arXiv:hep-th/9701081.
\end{thebibliography}